\documentclass[letterpaper,11pt]{moderncv}
\moderncvstyle{banking}
\moderncvcolor{black}

\usepackage[letterpaper, margin=0.8in]{geometry}
\usepackage{xcolor}
\usepackage{fontspec}
\setmainfont{Times New Roman}
\usepackage[unicode]{hyperref}

\name{Saeyoung}{Kim}
\address{Cambridge, MA, USA}{}{}
\phone[mobile]{+1 (351) 322 5510}
\email{saeyoung\_kim@uml.edu}
\social[linkedin]{saeyoung-macx-kim}

\recipient{Hiring Team}{Busek Co. Inc.\\Natick, MA}
\date{\today}
\opening{Dear Hiring Team,}
\closing{Sincerely,}

\begin{document}

\makelettertitle

I'm applying for the Manufacturing Engineer position. The description reads like a summary of what I've been doing for the past several years: hands-on assembly and testing of aerospace hardware, writing the procedures and documentation that make builds repeatable, and finding ways to make the process better each time.

At Hanwha Systems, I worked on military satellite programs where I did hands-on mechanical assembly and electrical integration in cleanroom environments. I wrote controlled test and assembly procedures that technicians used on the production floor, performed root cause analysis when builds had issues, and fed manufacturing problems back to the design teams. I also helped plan the facility architecture for a new high-throughput production line. The work required constant coordination across mechanical, electrical, software, and program teams, and everything had to be documented and traceable.

During my PhD, I took a CubeSat sensor payload through the full build cycle multiple times: circuit design, PCB fabrication, cleanroom MEMS processing, mechanical assembly, and system-level testing. Each revision, I refined the build process based on what broke or what was harder to assemble than it needed to be. That's where I developed the habit of treating documentation and process improvement as part of the engineering, not an afterthought.

More recently at UMass Lowell, I've been building Python tools for manufacturing process monitoring and anomaly detection, replacing commercial platforms with custom software that's easier to adapt. The work has reinforced my interest in systematizing production workflows and making them data-driven.

On the ITAR front: my green card is currently pending, and I expect to have it in hand within the next few months. I'm happy to discuss timing directly. I'm based in Cambridge, so Natick is a short commute.

\makeletterclosing

\end{document}
